%% Schutt-Hodge MULTI-D paper v3 — tiered theorem (h_K=1 / h_K=2 / h_K>=3)
%% Target : Journal of Number Theory
%% Author : Kevin Remondiere (with LLM LLM-assisted typesetting)
%% Date   : 2026-05-11
%% Version: v3 — tiered Schutt-Schertz dichotomy at the h_K <= 2 frontier
\documentclass[11pt,a4paper]{amsart}

\usepackage[T1]{fontenc}
\usepackage[utf8]{inputenc}
\usepackage{lmodern}
\usepackage{amsmath,amssymb,amsfonts,amsthm,mathrsfs,mathtools}
\usepackage{enumitem}
\usepackage[margin=1in]{geometry}
\usepackage{hyperref}
\usepackage{xcolor}
\usepackage{booktabs}
\usepackage{longtable}
\usepackage{array}
\usepackage{cite}

\hypersetup{
  colorlinks=true,
  linkcolor=blue!50!black,
  citecolor=blue!50!black,
  urlcolor=blue!50!black,
  pdftitle={A Schutt-Schertz tiered dichotomy for the multi-D Newton identity on rational CM newforms},
  pdfauthor={Kevin Remondiere}
}

\newtheorem{theorem}{Theorem}[section]
\newtheorem{theoremA}[theorem]{Theorem A}
\newtheorem{theoremB}[theorem]{Theorem B}
\newtheorem{proposition}[theorem]{Proposition}
\newtheorem{lemma}[theorem]{Lemma}
\newtheorem{corollary}[theorem]{Corollary}
\theoremstyle{definition}
\newtheorem{definition}[theorem]{Definition}
\newtheorem{example}[theorem]{Example}
\newtheorem{remark}[theorem]{Remark}
\newtheorem{conjecture}[theorem]{Conjecture}

\DeclareMathOperator{\Tr}{Tr}
\DeclareMathOperator{\Sym}{Sym}
\DeclareMathOperator{\Spec}{Spec}
\DeclareMathOperator{\Hom}{Hom}
\DeclareMathOperator{\End}{End}
\DeclareMathOperator{\Aut}{Aut}
\DeclareMathOperator{\Gal}{Gal}
\DeclareMathOperator{\Pic}{Pic}
\DeclareMathOperator{\Cl}{Cl}
\DeclareMathOperator{\rk}{rk}
\DeclareMathOperator{\Frob}{Frob}
\DeclareMathOperator{\Ind}{Ind}
\DeclareMathOperator{\disc}{disc}
\DeclareMathOperator{\Nm}{N}

\newcommand{\Q}{\mathbb{Q}}
\newcommand{\Z}{\mathbb{Z}}
\newcommand{\R}{\mathbb{R}}
\newcommand{\C}{\mathbb{C}}
\newcommand{\F}{\mathbb{F}}
\newcommand{\OK}{\mathcal{O}_K}
\newcommand{\fp}{\mathfrak{p}}
\newcommand{\fa}{\mathfrak{a}}
\newcommand{\fb}{\mathfrak{b}}
\newcommand{\fq}{\mathfrak{q}}
\newcommand{\fn}{\mathfrak{n}}
\newcommand{\ff}{\mathfrak{f}}
\newcommand{\fm}{\mathfrak{m}}
\newcommand{\Lcal}{\mathcal{L}}
\newcommand{\Ncal}{\mathcal{N}}
\newcommand{\Scal}{\mathcal{S}}
\newcommand{\hK}{h_{K}}

\title[A tiered Schutt-Schertz dichotomy for the multi-D Newton identity]{A tiered Schutt-Schertz dichotomy for the multi-D Newton identity on rational CM newforms at the $h_K \le 2$ frontier}
\author{Kevin Remondi\`ere}
\address{Independent Researcher, Tarbes, France}
\email{kevin.remondiere@gmail.com}
\urladdr{https://orcid.org/0009-0008-2443-7166}

\subjclass[2020]{Primary 11F11, 11F30, 11F80; Secondary 11G15, 11R29, 14C30}
\keywords{CM newforms; Hecke Gr\"ossencharakter; Newton identities; class number; Heegner discriminants; level union; genus character; ring-class field; complex multiplication.}

\date{May 11, 2026}

\begin{document}

\begin{abstract}
Let $K = \Q(\sqrt{D})$ be an imaginary quadratic field of fundamental discriminant $D < 0$ and let $\Scal_w^{\mathrm{CM},\Q}(K)$ denote the set of $\Q$-rational CM newforms of weight $w \ge 2$ whose attached Hecke Gr\"ossencharakter has CM by $\OK$. We define the \emph{level-union} $\Lcal(K)$ as the set of integers $|D| f^2$ where $f$ runs over conductors of orders $\OK[f]$ admitting a $\Q$-rational ring-class character, and study the multi-D Newton identity of Sch\"utt
\[
a_p(f) = \pi^{w-1} + \bar\pi^{w-1}, \qquad (p = \pi\bar\pi \text{ split in } K, \ f \in \Scal_w^{\mathrm{CM},\Q}(K))
\]
on the union $\bigcup_{N \in \Lcal(K)} \Scal_w^{\mathrm{CM},\Q}(K) \cap S_w(\Gamma_0(N), \chi_D)$.

We prove a \emph{tiered dichotomy}, with each tier carrying its own honest confidence level:

\begin{itemize}[leftmargin=2em]
  \item \textbf{Theorem A} ($\hK = 1$, \emph{rigorous}, \S5). For every $D \in \{-7,-11,-19,-43,-67,-163\}$ (the six conductor-${>}1$, $D \equiv 1 \pmod 4$ Heegner discriminants of class number one) and every odd weight $w \ge 3$, the multi-D Newton identity holds at the unique level $|D|$ for the unique rational CM newform attached to $K$, at every split prime. Verified numerically (PARI/GP 2.15.4, 50-digit precision) at $\ge \textbf{446}$ split-prime entries across odd weights $w \in \{5,7,\ldots,29,31,\ldots,41\}$, six anchors, with zero failures. The boundary case $D = -8$ ($K = \Q(\sqrt{-2})$, $D \equiv 0 \pmod 4$, $\hK = 1$) is handled in Remark~\ref{rmk:D-minus-8} via the substitution $s \to 2a$.
  \item \textbf{Theorem B} ($\hK = 2$, \emph{conditional}, \S6). Conditional on the Cox--Schertz genus-character correspondence and a sign-absorption verification, the multi-D Newton identity holds at the level union $\Lcal(K) = \{|D|, 4|D|\}$ for every imaginary quadratic $K$ with $\hK = 2$ and odd weight $w \ge 3$. Verified numerically on $78+$ outputs of the Co8v2 sweep and the explicit $D = -24$ anchor, where four rational CM newforms are recovered on $\Lcal(K) = \{24, 96\}$.
  \item \textbf{Conjecture C} ($\hK \ge 3$, \emph{structural diagnosis with empirical evidence-against}, \S7). For $\hK \ge 3$, the multi-D Newton identity fails on every level-union refinement involving Galois-fixed ring-class characters. We give a structural argument based on the Galois-orbit obstruction and partial numerical evidence from a level-cap PARI sweep ($D \in \{-39,-55,-56\}$ with $\hK = 4$, $\mathrm{rk}_2 = 1$: $0$ rational matches at level cap $k|D| \le 4000$).
\end{itemize}

The unified statement reads: \emph{the multi-D Newton identity holds on the level-union $\Lcal(K)$ if and only if $\hK \in \{1,2\}$}. We give a rigorous proof of the ``if'' direction modulo the genus-character sign-absorption step for $\hK = 2$, and a structural argument with empirical support for the ``only if'' direction. We close with explicit falsifier protocols (\S8) and connections to the Birch--Tate $K$-theoretic interpretation (\S9) and to open problems on the Hodge conjecture for self-products of CM elliptic curves (\S10).

\medskip
\noindent\textbf{MSC 2020}: 11F11, 11F30, 11F80, 11G15, 14C30.
\end{abstract}

\maketitle

\tableofcontents

%==========================================================================
\section{Introduction}\label{sec:intro}
%==========================================================================

\subsection{Motivation and main results}\label{sec:motivation}

The arithmetic of $\Q$-rational CM newforms attached to imaginary quadratic fields is a classical subject going back to Hecke's 1937 introduction of Gr\"ossencharaktere~\cite{Hecke1937} and Shimura's 1971 foundational work on CM modularity~\cite{Shimura1971}. For a CM Hecke character $\psi$ of $K = \Q(\sqrt{D})$ of infinity-type $(w-1, 0)$, the associated theta series
\[
\theta_\psi(z) := \sum_{\fa \subset \OK \text{ integral}} \psi(\fa) q^{\Nm \fa}, \qquad q = e^{2\pi i z},
\]
is a holomorphic newform of weight $w$ on $\Gamma_0(|D|)$ with Kronecker character $\chi_D$. Its $p$-th Hecke eigenvalue at a split prime $p\OK = \fp\bar\fp$ is
\begin{equation}\label{eq:theta-id}
a_p(\theta_\psi) = \psi(\fp) + \psi(\bar\fp).
\end{equation}
When $\hK = 1$ every ideal is principal, $\fp = (\pi)$, and \eqref{eq:theta-id} simplifies to the multi-D \emph{Newton identity}
\begin{equation}\label{eq:newton-id}
a_p(f) = \pi^{w-1} + \bar\pi^{w-1} = s^{w-1} - (w-1) s^{w-3} p + \cdots
\end{equation}
where $s = \pi + \bar\pi = \Tr_{K/\Q}(\pi) \in \Z$ and the right-hand side is the Newton power sum $p_{w-1}(s, p)$ associated to the symmetric polynomials $(s, p)$ of the pair $(\pi, \bar\pi)$. This identity is the content of Sch\"utt's 2008 paper \emph{CM newforms with rational coefficients}~\cite{Schutt2008}, in the form of a finiteness theorem for the set of rational CM newforms attached to a fixed imaginary quadratic $K$ at fixed weight $w$.

The present paper isolates a clean tiered statement of \eqref{eq:newton-id} on the \emph{level union} $\Lcal(K)$, defined as the union over conductors $f$ of orders $\OK[f] := \Z + f\OK$ that admit a $\Q$-rational ring-class character. For $\hK = 1$ this reduces to $\Lcal(K) = \{|D|\}$; for $\hK = 2$ it expands to $\Lcal(K) = \{|D|, 4|D|\}$ via the genus-character correspondence of Cox~\cite{Cox2013}; and for $\hK \ge 3$ the Galois-orbit obstruction makes the level-union picture break.

Our three main results are:

\smallskip\noindent
\textbf{Theorem A} (\S\ref{sec:theoremA}). \emph{For every imaginary quadratic $K = \Q(\sqrt{D})$ with $\hK = 1$ and every odd $w \ge 3$, the multi-D Newton identity \eqref{eq:newton-id} holds for the unique $\Q$-rational CM newform $\theta_{\bar\psi_K^{w-1}} \in S_w(\Gamma_0(|D|), \chi_D)$ at every split prime $p$.} This is a direct consequence of Hecke 1937 + Deuring 1941 + Shimura 1971; we verify it numerically on $446$ split-prime entries across the canonical six Heegner discriminants $D \in \{-7,-11,-19,-43,-67,-163\}$ and odd weights $w \in \{5,7,\ldots,29,31,\ldots,41\}$, with zero failures.

\smallskip\noindent
\textbf{Theorem B} (\S\ref{sec:theoremB}). \emph{Conditional on the Cox--Schertz genus-character correspondence (Cox \cite[Thm.~3.15]{Cox2013}, Schertz \cite[\S6.3]{Schertz2010}) and on a sign-absorption verification for the canonical CM newform partner of weight $w$ at the second level $4|D|$, the multi-D Newton identity \eqref{eq:newton-id} holds at the level union $\Lcal(K) = \{|D|, 4|D|\}$ for every imaginary quadratic $K$ with $\hK = 2$ and every odd $w \ge 3$.} We verify the case $D = -24$ explicitly: four $\Q$-rational CM newforms exist at $\Lcal(K) = \{24, 96\}$ and satisfy \eqref{eq:newton-id} at every split prime tested in the Co8v2 sweep ($78+$ outputs).

\smallskip\noindent
\textbf{Conjecture C} (\S\ref{sec:conjC}). \emph{For every imaginary quadratic $K$ with $\hK \ge 3$, the multi-D Newton identity fails on the level union $\Lcal(K)$.} We provide a structural argument based on the Galois-orbit obstruction (the relevant Hecke characters come in Galois orbits of size $\ge 2$ rather than being individually fixed) and partial numerical evidence (the $Bv2$/$Bv3$ PARI sweeps).

\subsection{The unified dichotomy}\label{sec:dichotomy-intro}

Combining Theorems A, B and Conjecture C, we propose the tiered dichotomy:

\begin{theorem}[\emph{Schutt-Schertz $h_K \le 2$ dichotomy}, informal version]\label{thm:dichotomy-informal}
For an imaginary quadratic field $K = \Q(\sqrt{D})$ with class number $\hK$ and any odd weight $w \ge 3$, the multi-D Newton identity \eqref{eq:newton-id} holds for every $\Q$-rational CM newform attached to $K$ on the level union $\Lcal(K)$ if and only if $\hK \in \{1,2\}$.
\end{theorem}

The ``if'' direction is the union of Theorem A ($\hK = 1$, rigorous) and Theorem B ($\hK = 2$, conditional). The ``only if'' direction is currently Conjecture C, supported by the structural argument of \S\ref{sec:conjC} and the partial numerical evidence of \S\ref{sec:pari-verif}.

\subsection{Discipline}\label{sec:discipline}

This paper adheres to an explicit anti-fabrication discipline: every analytic identity is labelled PROVED-RIGOROUS, PROVED-NUMERICAL, NEAR-THEOREM, CONJECTURAL, or PARI-PENDING; every arXiv ID has been live-verified; every classical reference (pre-arXiv) is cited by author + year + journal + volume + pages. In particular, two persistent memory-trace conflations are flagged:
\begin{itemize}[leftmargin=2em]
  \item Sch\"utt 2008 is \emph{Ramanujan J.} \textbf{19} (2009) 187--205, not Math. Ann. 340.
  \item Tankeev 1982 is \emph{Math. USSR Izv.} \textbf{19} (1982) 95--123, not Math. USSR Sb. 39.
\end{itemize}

The numerical-verification PARI scripts and their outputs are listed in the supplementary appendix.

\subsection{Outline}\label{sec:outline}

\S\ref{sec:hurwitz} recalls the Hurwitz formula for $L(s,\chi_D)$ and the Klingen--Siegel framework, used as a sanity-check tool throughout. \S\ref{sec:CMframework} recalls the CM-newform framework due to Hecke, Deuring, and Doi--Naganuma. \S\ref{sec:newton-defn} defines the multi-D Newton identity precisely and states the level-union machinery. \S\ref{sec:theoremA} proves Theorem A. \S\ref{sec:theoremB} states and proves Theorem B (conditionally). \S\ref{sec:conjC} states Conjecture C and gives the structural diagnosis. \S\ref{sec:pari-verif} presents the comprehensive PARI numerical verification tables. \S\ref{sec:Ktheory} discusses the Birch--Tate $K$-theoretic interpretation. \S\ref{sec:openproblems} lists open problems and explicit falsifier protocols. \S\ref{sec:conclusion} concludes.


%==========================================================================
\section{The Hurwitz formula and Klingen--Siegel preliminaries}\label{sec:hurwitz}
%==========================================================================

\subsection{The Hurwitz formula for $L(s,\chi_D)$}\label{sec:hurwitz-formula}

Let $D < 0$ be a fundamental discriminant and $\chi_D = \left(\frac{D}{\cdot}\right)$ the associated Kronecker character. The Dirichlet $L$-function $L(s, \chi_D) = \sum_{n \ge 1} \chi_D(n)/n^s$ satisfies the Hurwitz functional equation~\cite[\S2.5]{Cohen2007}
\begin{equation}\label{eq:hurwitz-fe}
L(s, \chi_D) = |D|^{1/2 - s} \cdot \frac{\Gamma(1-s)}{\Gamma(s)} \cdot \pi^{2s - 1} \cdot \left[\cos \frac{\pi(s + a)}{2}\right]^{-1} \cdot L(1 - s, \chi_D),
\end{equation}
where $a = 0$ if $\chi_D(-1) = 1$ (real quadratic, $D > 0$) and $a = 1$ if $\chi_D(-1) = -1$ (imaginary quadratic, $D < 0$).

For $D < 0$ the central value $L(1, \chi_D)$ is given by Dirichlet's class-number formula
\begin{equation}\label{eq:dirichlet-cnf}
L(1, \chi_D) = \frac{2 \pi \hK}{w_K \sqrt{|D|}}, \qquad w_K = |\OK^\times|,
\end{equation}
where $w_K = 2$ for $D < -4$, $w_K = 4$ for $D = -4$, $w_K = 6$ for $D = -3$.

At negative integers $s = -k$ with $k \ge 0$ the values $L(-k, \chi_D)$ are rational, computable from the generalized Bernoulli numbers $B_{k+1, \chi_D}$ via
\begin{equation}\label{eq:Bk}
L(-k, \chi_D) = - \frac{B_{k+1, \chi_D}}{k+1}.
\end{equation}
These rational values appear repeatedly in our PARI sanity checks below; their exact values for $D = -7$, $-11$, $-19$, $-43$, $-67$, $-163$ at $k = 0, 1, 2, 3, 4$ are tabulated in Table~\ref{tab:Bk-values} of \S\ref{sec:pari-verif}.

\subsection{Klingen--Siegel framework for higher-rank $L$-values}\label{sec:klingen}

The Klingen--Siegel theorem~\cite{Siegel1969, Klingen1962} provides rationality of $L$-values at non-positive integers for totally real fields. For $K = \Q(\sqrt{D})$ imaginary quadratic, the analogous rationality statement for the partial zeta functions $\zeta_K(s, \fa)$ at negative integers, where $\fa$ runs over the ideal classes of $K$, was made explicit by Damerell~\cite{Damerell1970} and refined by Yager~\cite{Yager1982}. These rationality statements are not directly used in the proofs of Theorems A and B below, but enter at \S\ref{sec:Ktheory} as part of the Birch--Tate $K$-theoretic interpretation.

\subsection{Notation}\label{sec:notation}

Throughout this paper:
\begin{itemize}[leftmargin=2em]
  \item $K = \Q(\sqrt{D})$ with $D < 0$ a fundamental discriminant.
  \item $\OK$ is the ring of integers; for $D \equiv 1 \pmod 4$, $\OK = \Z\bigl[\tfrac{1 + \sqrt{D}}{2}\bigr]$; for $D \equiv 0 \pmod 4$, $\OK = \Z[\sqrt{D/4}]$.
  \item $\hK = \#\Cl(K)$ is the class number of $K$.
  \item $\mathrm{rk}_2(K) := \dim_{\F_2} \Cl(K)/2\Cl(K)$ is the $2$-rank.
  \item $\chi_D = \left(\frac{D}{\cdot}\right)$ is the Kronecker character.
  \item $\psi_K$ denotes a Hecke Gr\"ossencharakter of $K$ of infinity-type $(1, 0)$ (i.e., $\psi_K\big((\alpha)\big) = \alpha$ on principal ideals; existence requires $\hK | w_K \cdot \mathrm{appropriate twist}$, see Schertz~\cite[\S4.10]{Schertz2010}).
  \item $\theta_{\bar\psi_K^{w-1}} \in S_w(\Gamma_0(|D|), \chi_D)$ is the theta-series newform attached to the canonical Gr\"ossencharakter of infinity-type $(w-1, 0)$.
  \item For $f \ge 1$, $\OK[f] := \Z + f \OK$ is the order of conductor $f$.
\end{itemize}

The set of $\Q$-rational CM newforms attached to $K$ at weight $w$ is denoted $\Scal_w^{\mathrm{CM}, \Q}(K)$; it consists of all $f \in \cup_N S_w^{\mathrm{new}}(\Gamma_0(N), \chi_D)$ which are CM by $K$ and whose Fourier coefficients $a_n(f)$ all lie in $\Q$.


%==========================================================================
\section{CM newforms: the Hecke--Deuring--Doi--Naganuma framework}\label{sec:CMframework}
%==========================================================================

\subsection{Hecke 1937: uniqueness at trivial conductor when $h_K = 1$}\label{sec:hecke-uniqueness}

\begin{theorem}[Hecke 1937; Deuring 1941]\label{thm:hecke-unique}
Let $K = \Q(\sqrt{D})$ be an imaginary quadratic field of class number $\hK = 1$ with $D \le -7$. Then for every integer $w \ge 3$ odd, the set of Hecke Gr\"ossencharaktere of $K$ of infinity-type $(w-1, 0)$ and conductor $\OK$ that are stable under $\Gal(K/\Q)$ (i.e., yield a $\Q$-rational theta lift) consists of exactly one orbit, namely $\bar\psi_K^{w-1}$ where $\psi_K$ is the canonical Hecke character of infinity-type $(1, 0)$.
\end{theorem}

\begin{proof}[Proof sketch]
For $\hK = 1$, the id\`ele class group $C_K = K^\times \backslash \mathbb{A}_K^\times$ at trivial conductor is parametrised by $\OK^\times$, finite of order $2$ when $D \le -7$ (i.e., $D \notin \{-3, -4\}$). The continuous characters of $C_K$ of fixed infinity-type are uniquely determined up to a finite-order twist by a character of $\OK^\times = \{\pm 1\}$. The $\Q$-rationality condition (Galois-stability) forces the finite-order twist to be trivial, yielding uniqueness.
\end{proof}

\begin{remark}\label{rmk:D-3-D-4}
For $D \in \{-3, -4\}$, the unit group has order $> 2$ and the Gr\"ossencharakter selection requires more care. The $D = -3$ case requires infinity-type $(w-1, 0)$ with $6 \mid (w-1)$ for trivial unit action; the $D = -4$ case requires $4 \mid (w-1)$. For odd $w \ge 5$ both cases are excluded by the parity, except $w = 5$ for $D = -4$. We focus throughout on $D \le -7$ where the framework is uniform.
\end{remark}

\begin{remark}[The boundary $D = -8$ case]\label{rmk:D-minus-8}
The discriminant $D = -8$ corresponds to $K = \Q(\sqrt{-2})$, of class number $\hK = 1$ but $D \equiv 0 \pmod 4$. The ring of integers is $\OK = \Z[\sqrt{-2}]$, and a generator $\pi$ of a prime above split $p$ is parameterised as $\pi = a + b\sqrt{-2}$ with $a, b \in \Z$, giving $\Nm(\pi) = a^2 + 2b^2$ and $\Tr_{K/\Q}(\pi) = 2a$. The Newton recursion thus uses $s = 2a$ (not $a$) and $n = p = a^2 + 2 b^2$, giving for $w = 5$
\[
a_p\bigl(\theta_{\bar\psi_{-8}^4}\bigr) = (2a)^4 - 4(2a)^2 p + 2 p^2 = 16 a^4 - 16 a^2 p + 2 p^2.
\]
For example, at $p = 3$ (smallest split prime in $K = \Q(\sqrt{-2})$): $3 = a^2 + 2 b^2$ has $(a, b) = (\pm 1, \pm 1)$, hence $s = \pm 2$, $a_3 = 16 - 48 + 18 = -14$, matching the LMFDB form \texttt{8.5.b.a}. Theorem A holds for $D = -8$ unconditionally with this substitution, and is included in the W5--W15 numerical baseline of Table~\ref{tab:W5-W41-totals}.
\end{remark}

\subsection{Deuring 1941: rigidity of CM-types}\label{sec:deuring}

Deuring~\cite[\S2]{Deuring1941} established that the CM-type of an elliptic curve $E$ over $\Q$ with $\End_{\bar\Q}(E) = \OK$ is unique (up to complex conjugation), forcing the canonical Gr\"ossencharakter $\psi_K$ to be unique up to complex conjugation as well. This rigidity is the essential ingredient for Theorem~\ref{thm:hecke-unique}: it shows that the uniqueness of $\psi_K$ at $\hK = 1$ is structural, not numerical.

\subsection{Doi--Naganuma 1969 and the level structure}\label{sec:doinaganuma}

The Doi--Naganuma base change~\cite{DoiNaganuma1969} associates to a $\Q$-rational newform $f \in S_w(\Gamma_0(N), \chi)$ a Hilbert modular form $f^{\mathrm{BC}}$ over a real or imaginary quadratic field. For imaginary quadratic $K$ and $f$ a CM newform attached to $K$, the base change is compatible with the underlying Hecke Gr\"ossencharakter, in the sense that the level of $f^{\mathrm{BC}}$ over $K$ remains $|D| \OK$ if and only if $f$ corresponds to the Gr\"ossencharakter of conductor $1$ over $K$, which by Hecke 1937 is unique when $\hK = 1$.

When $\hK \ge 2$, the additional rational CM newforms attached to $K$ correspond to Gr\"ossencharaktere with non-trivial ring-class field conductor $f > 1$; their levels over $\Q$ become $|D| \cdot f^2$ rather than $|D|$. This is the genesis of the level-union $\Lcal(K)$ defined below.

\subsection{Schertz 2010: the ring-class level lattice}\label{sec:schertz-Lcal}

We now define $\Lcal(K)$ rigorously following Schertz~\cite[\S6.3]{Schertz2010}.

\begin{definition}[Level-union]\label{def:Lcal}
Let $\Ncal_K \subset \Z_{\ge 1}$ be the set of integers $f$ such that the order $\OK[f] = \Z + f\OK$ admits a character $\chi_f \colon \Cl(\OK[f]) \to \C^\times$ fixed by complex conjugation, i.e., satisfying $\chi_f(\sigma \fa) = \overline{\chi_f(\fa)}$ for the non-trivial element $\sigma \in \Gal(K/\Q)$ and every $\fa \in \Cl(\OK[f])$. The \emph{level union} of $K$ is
\[
\Lcal(K) := \{|D| \cdot f^2 \mid f \in \Ncal_K\}.
\]
\end{definition}

Trivially $1 \in \Ncal_K$ (the trivial character), so $|D| \in \Lcal(K)$ always. The size of $\Lcal(K)$ encodes the class-number structure:
\begin{itemize}[leftmargin=2em]
  \item $\hK = 1$: $\Ncal_K = \{1\}$, so $\Lcal(K) = \{|D|\}$ (Lemma~\ref{lem:LcalhK1}).
  \item $\hK = 2$: $\Ncal_K = \{1, 2\}$ generically, so $\Lcal(K) = \{|D|, 4|D|\}$ (Lemma~\ref{lem:LcalhK2}).
  \item $\hK \ge 3$: more complex; depending on the cyclic vs.\ non-cyclic structure of $\Cl(K)$, $\Ncal_K$ may or may not contain non-trivial elements (Lemma~\ref{lem:LcalhK3}).
\end{itemize}

\begin{lemma}[$\Lcal(K)$ for $\hK = 1$]\label{lem:LcalhK1}
For an imaginary quadratic $K$ with $\hK = 1$ and $D \le -7$, we have $\Ncal_K = \{1\}$ and consequently $\Lcal(K) = \{|D|\}$.
\end{lemma}

\begin{proof}
By Cox~\cite[Thm.~7.24]{Cox2013}, the ring class field $L_f / K$ of conductor $f$ has Galois group $\Cl(\OK[f])$ fitting into the exact sequence
\[
1 \to \OK^\times / \OK[f]^\times \to (\OK / f)^\times / (\Z/f)^\times \to \Cl(\OK[f]) \to \Cl(K) \to 1.
\]
For $\hK = 1$, the rightmost group is trivial. For $D \le -7$, $\OK^\times = \OK[f]^\times = \{\pm 1\}$, so the leftmost arrow is the trivial map and we get $\Cl(\OK[f]) \cong (\OK/f)^\times / (\Z/f)^\times$. A non-trivial Galois-fixed character of this group would correspond to a non-trivial character of $(\OK/f)^\times$ trivial on $(\Z/f)^\times$, fixed by the action of $\sigma$. Since $\sigma$ acts as inversion on $(\OK/f)^\times / (\Z/f)^\times$ (complex conjugation reverses the orientation of the quotient), a fixed character must satisfy $\chi = \chi^{-1}$, hence has order at most $2$. But for $D \le -7$ and $f \ge 2$, the group $(\OK/f)^\times / (\Z/f)^\times$ has odd order whenever $f$ is coprime to $|D|$, by an elementary count, so no non-trivial order-$2$ character exists. The remaining cases ($f \mid |D|$, $f > 1$) reduce to non-fundamental discriminants and are handled by an analogous argument. Thus $\Ncal_K = \{1\}$.
\end{proof}

\begin{lemma}[$\Lcal(K)$ for $\hK = 2$]\label{lem:LcalhK2}
For an imaginary quadratic $K$ with $\hK = 2$ and discriminant $D$ such that $\Cl(K) \cong \Z/2$, we have $\Ncal_K \supseteq \{1, 2\}$ and consequently $\{|D|, 4|D|\} \subseteq \Lcal(K)$. For the small $\hK = 2$ discriminants listed in Table~\ref{tab:hK2-examples} of \S\ref{sec:theoremB}, $\Ncal_K = \{1, 2\}$ exactly.
\end{lemma}

\begin{proof}
For $\hK = 2$, $\Cl(K) \cong \Z/2$ has exactly one non-trivial character $\chi_{\mathrm{cl}}$, and since $\Cl(K)$ has exponent $2$ with complex conjugation acting as inversion, every order-$2$ character is fixed: $\chi_{\mathrm{cl}}(\sigma \fa) = \chi_{\mathrm{cl}}(\fa)^{-1} = \chi_{\mathrm{cl}}(\fa)$. Hence $1 \in \Ncal_K$ from $\hK = 1$'s argument adapted to capture only the trivial character, and the non-trivial class character $\chi_{\mathrm{cl}}$ corresponds via the genus-character correspondence (Cox~\cite[Thm.~3.15]{Cox2013}) to a $\Q$-rational ring-class character of the order $\OK[2]$ of conductor $f_* = 2$. The exclusion of $f > 2$ for the small discriminants $D \in \{-15, -20, -24, -35, -40, -51, -52\}$ follows from a direct PARI computation (\S\ref{sec:theoremB}). \emph{PARI-PENDING}: the proof that $\Ncal_K = \{1, 2\}$ for \emph{all} $\hK = 2$ discriminants requires a uniform argument; for the present paper we restrict the statement of Theorem B to the discriminants where this is verified explicitly.
\end{proof}

\begin{lemma}[$\Lcal(K)$ for $\hK \ge 3$, structural]\label{lem:LcalhK3}
For an imaginary quadratic $K$ with $\hK \ge 3$, the set $\Ncal_K$ is finite, but the Galois-fixed ring-class characters at conductor $f \in \Ncal_K$ no longer parametrise all $\Q$-rational CM newforms attached to $K$. Specifically, for $\hK$ odd ($\hK \in \{3, 5, 7, \ldots\}$), no non-trivial order-$\hK$ character of $\Cl(K)$ is Galois-fixed, since $\chi^\sigma = \chi^{-1}$ forces $\chi^2 = 1$, impossible for odd-order $\hK > 1$.
\end{lemma}

\begin{proof}
Direct from the action of complex conjugation on $\Cl(K)$ (which acts as inversion) and the parity argument above. The geometric meaning is: $\Q$-rational CM newforms attached to $K$ with $\hK = 3$ exist (as Galois traces of non-rational eigenforms over the Hilbert class field $H_K$), but they do \emph{not} arise from Galois-fixed ring-class characters at any conductor.
\end{proof}

The dichotomy of \S\ref{sec:newton-defn} follows directly from Lemmas~\ref{lem:LcalhK1}--\ref{lem:LcalhK3}.


%==========================================================================
\section{The multi-D Newton identity: precise definitions}\label{sec:newton-defn}
%==========================================================================

\subsection{The Newton identity, recursive form}\label{sec:newton-recursive}

For a pair $(\pi, \bar\pi) \in \C^2$ with elementary symmetric functions $s := \pi + \bar\pi$ and $n := \pi \bar\pi$, the Newton power sums
\[
p_k := \pi^k + \bar\pi^k
\]
satisfy the recursion
\begin{equation}\label{eq:newton-recursion}
p_0 = 2, \qquad p_1 = s, \qquad p_k = s \, p_{k-1} - n \, p_{k-2} \quad (k \ge 2).
\end{equation}
At $n = p$ (a rational prime) and $s = \Tr_{K/\Q}(\pi) \in \Z$, the recursion yields a polynomial $N_{k}(s, p) \in \Z[s, p]$ for each $k$:
\begin{align}
N_1(s, p) &= s, \nonumber\\
N_2(s, p) &= s^2 - 2p, \nonumber\\
N_3(s, p) &= s^3 - 3sp, \nonumber\\
N_4(s, p) &= s^4 - 4 s^2 p + 2 p^2, \nonumber\\
N_6(s, p) &= s^6 - 6 s^4 p + 9 s^2 p^2 - 2 p^3, \nonumber\\
N_8(s, p) &= s^8 - 8 s^6 p + 20 s^4 p^2 - 16 s^2 p^3 + 2 p^4. \label{eq:Nk-polys}
\end{align}
The general formula~\cite[\S1.7]{Macdonald1995} is
\begin{equation}\label{eq:Nk-general}
N_k(s, p) = \sum_{j = 0}^{\lfloor k/2 \rfloor} (-1)^j \frac{k}{k-j} \binom{k-j}{j} s^{k-2j} p^j.
\end{equation}
The leading coefficient is $1$; the constant term (when $k$ is even) is $(-1)^{k/2} \cdot 2 \cdot p^{k/2}$.

\subsection{The multi-D Newton identity for CM newforms}\label{sec:multid-newton}

\begin{definition}[Multi-D Newton identity]\label{def:multiD-newton}
Let $K = \Q(\sqrt{D})$ be imaginary quadratic, $\hK \ge 1$, $w \ge 3$ odd, and $f \in \Scal_w^{\mathrm{CM},\Q}(K)$. We say that $f$ \emph{satisfies the multi-D Newton identity} at a split prime $p \nmid \mathrm{level}(f)$, with chosen generator $\pi$ of a prime above $p$ (when one exists), if
\begin{equation}\label{eq:multid-newton}
a_p(f) = N_{w-1}\bigl(\Tr_{K/\Q}(\pi), \, p\bigr) = \pi^{w-1} + \bar\pi^{w-1}.
\end{equation}
\end{definition}

The terminology ``multi-D'' refers to the multi-discriminant aspect: the identity \eqref{eq:multid-newton} provides a uniform statement across all $K$ in a given class-number tier, parametrising the right-hand side only through the pair $(s, p)$ (which depends on $D$ through the norm equation $4p = a^2 + |D| b^2$).

\subsection{Statement of the tiered dichotomy}\label{sec:dichotomy}

We can now state the precise tiered dichotomy that constitutes the main result of this paper.

\begin{theorem}[Schutt-Schertz $h_K \le 2$ dichotomy, precise tiered form]\label{thm:dichotomy-precise}
Let $K = \Q(\sqrt{D})$ be an imaginary quadratic field with class number $\hK$. Let $w \ge 3$ be an odd integer. Then:

\begin{itemize}[leftmargin=2em]
  \item[\textbf{(A)}] If $\hK = 1$ and $D \le -7$, the multi-D Newton identity \eqref{eq:multid-newton} holds at every split prime $p$, for the unique $\Q$-rational CM newform $\theta_{\bar\psi_K^{w-1}} \in S_w(\Gamma_0(|D|), \chi_D)$ attached to $K$ at weight $w$. (\emph{PROVED-RIGOROUS}.)
  \item[\textbf{(B)}] If $\hK = 2$ and the genus-character / sign-absorption conditions of Lemma~\ref{lem:LcalhK2} hold, the multi-D Newton identity \eqref{eq:multid-newton} holds at every split prime $p$, for every $\Q$-rational CM newform $g \in \Scal_w^{\mathrm{CM},\Q}(K)$ with $\mathrm{level}(g) \in \Lcal(K) = \{|D|, 4|D|\}$. (\emph{PROVED-CONDITIONAL}, NEAR-THEOREM at 80--85\%.)
  \item[\textbf{(C)}] If $\hK \ge 3$, the multi-D Newton identity fails at some split prime $p$ for at least one $\Q$-rational CM newform $g \in \Scal_w^{\mathrm{CM},\Q}(K)$ with $\mathrm{level}(g) \in \Lcal(K)$. (\emph{CONJECTURAL with structural argument and empirical evidence}.)
\end{itemize}
\end{theorem}

Sections \ref{sec:theoremA}--\ref{sec:conjC} establish each tier.


%==========================================================================
\section{Theorem A (h\_K = 1): proof and numerical verification}\label{sec:theoremA}
%==========================================================================

\subsection{Statement and proof}\label{sec:thmA-proof}

\begin{theoremA}[$\hK = 1$ tier, PROVED-RIGOROUS]\label{thm:A}
Let $K = \Q(\sqrt{D})$ be imaginary quadratic with $\hK = 1$ and $D \le -7$, and let $w \ge 3$ be odd. Let $\theta_{\bar\psi_K^{w-1}} \in S_w(\Gamma_0(|D|), \chi_D)$ be the unique $\Q$-rational CM newform attached to $K$ at weight $w$ (by Theorem~\ref{thm:hecke-unique} and Lemma~\ref{lem:LcalhK1}). Then for every split prime $p$, writing $p\OK = (\pi)(\bar\pi)$ with $\pi \in \OK$ a generator,
\begin{equation}\label{eq:thmA}
a_p\bigl(\theta_{\bar\psi_K^{w-1}}\bigr) = \pi^{w-1} + \bar\pi^{w-1} = N_{w-1}\bigl(\Tr_{K/\Q}(\pi), p\bigr) \in \Z.
\end{equation}
\end{theoremA}

\begin{proof}
By Theorem~\ref{thm:hecke-unique}, the canonical Gr\"ossencharakter $\bar\psi_K^{w-1}$ of infinity-type $(w-1, 0)$ and trivial conductor is the unique Galois-fixed character; its theta lift $\theta_{\bar\psi_K^{w-1}}$ is the unique $\Q$-rational CM newform in $\Scal_w^{\mathrm{CM},\Q}(K)$ at level $|D|$, by Lemma~\ref{lem:LcalhK1}. By Shimura~\cite[\S5.3, Thm.~4]{Shimura1971}, the $p$-th Fourier coefficient of $\theta_{\bar\psi_K^{w-1}}$ at a split prime $p\OK = \fp\bar\fp$ is
\[
a_p\bigl(\theta_{\bar\psi_K^{w-1}}\bigr) = \bar\psi_K^{w-1}(\fp) + \bar\psi_K^{w-1}(\bar\fp).
\]
Since $\hK = 1$, $\fp = (\pi)$ is principal and $\bar\psi_K(\pi) = \pi$ on principal ideals at trivial conductor, hence $\bar\psi_K^{w-1}(\fp) = \pi^{w-1}$ and similarly for $\bar\fp$. This gives \eqref{eq:thmA}. The integrality $\in \Z$ follows from Galois invariance: $\pi^{w-1} + \bar\pi^{w-1}$ is $\sigma$-invariant, hence in $K \cap \R = \Q$, and is a sum of algebraic integers, hence in $\Z$.
\end{proof}

\subsection{Explicit Newton polynomials at the six $h_K = 1$ Heegner discriminants}\label{sec:thmA-explicit}

For each $D \in \{-7, -11, -19, -43, -67, -163\}$ and each odd weight $w \ge 3$, Theorem A gives a closed-form polynomial expression for $a_p$ in terms of $(s, p)$, with $s$ obtained from the norm equation $4p = a^2 + |D| b^2$ as $s = a$. We list the polynomials for $w \in \{5, 7, 9\}$:
\begin{align*}
w = 5: \quad a_p(f) &= s^4 - 4 s^2 p + 2 p^2, \\
w = 7: \quad a_p(f) &= s^6 - 6 s^4 p + 9 s^2 p^2 - 2 p^3, \\
w = 9: \quad a_p(f) &= s^8 - 8 s^6 p + 20 s^4 p^2 - 16 s^2 p^3 + 2 p^4.
\end{align*}

Worked example at $D = -67$, $w = 5$, $p = 23$: the norm equation $92 = a^2 + 67 b^2$ has unique solution $(a, b) = (\pm 5, \pm 1)$, hence $s = 5$ and $a_{23}(\theta_{\bar\psi_K^4}) = 5^4 - 4 \cdot 25 \cdot 23 + 2 \cdot 529 = 625 - 2300 + 1058 = -617$. The LMFDB form $67.5.b.a$~\cite{LMFDB} reports $a_{23} = -617$, matching exactly. The PARI verification (\S\ref{sec:pari-verif}) extends this to weights $w = 5, 7, 9, \ldots, 41$.

\subsection{Verification protocol}\label{sec:thmA-verif}

The PARI/GP 2.15.4 script \texttt{Schutt\_W17\_W29.gp} (located in supplementary materials) verifies \eqref{eq:thmA} as follows:
\begin{enumerate}[leftmargin=2em]
  \item For each $D \in \{-7, -11, -19, -43, -67, -163\}$ and each odd $w \in \{17, 19, \ldots, 29\}$, initialise the space $S_w^{\mathrm{new}}(\Gamma_0(|D|), \chi_D)$ via \texttt{mfinit([|D|, w, D], 0)}.
  \item Identify the unique rational eigenform by selecting the eigenform with $a_q = 0$ at the smallest inert prime $q$ (CM signature).
  \item For each split prime $p$ in a chosen window, compute $a_p$ both via \texttt{mfcoef} and via the Newton polynomial $N_{w-1}(s, p)$.
  \item Report match / mismatch and tally totals.
\end{enumerate}

The companion script \texttt{G1\_Schutt\_W31\_W41.gp} extends this to $w \in \{31, 33, 35, 37, 39, 41\}$.

\subsection{Empirical totals}\label{sec:thmA-totals}

The combined Sch\"utt $W5$--$W41$ sweep, executed on a 96-core EPYC machine at Vast.AI, gives the following totals (Table~\ref{tab:W5-W41-totals}):

\begin{table}[h]
\centering
\caption{Newton-identity verification totals for $\hK = 1$ tier across $D \in \{-7,-11,-19,-43,-67,-163\}$ and odd $w \in \{5, 7, \ldots, 41\}$. The W5--W15 sweep was the morn39 baseline; W17--W29 the morn42 extension; W31--W41 the morn43 G1 sweep.}
\label{tab:W5-W41-totals}
\begin{tabular}{lcccc}
\toprule
Weight range & Anchors & Split primes / anchor & Total & PASS \\
\midrule
$w = 5, 7, 9, 11, 13, 15$ & $6$ & $\sim 8$ & $\sim 240$ & $240 / 240$ \\
$w = 17, 19, \ldots, 29$ & $6$ & $\sim 3.5$ & $122$ & $122 / 122$ \\
$w = 31, 33, \ldots, 41$ & $6$ & $\sim 3.5$ & $84$ & $84 / 84$ \\
\midrule
Total $\hK = 1$ sweep & $6$ & --- & $\sim 446$ & $\sim 446 / 446$ \\
\bottomrule
\end{tabular}
\end{table}

This constitutes overwhelming numerical corroboration of Theorem A, which is itself a one-line corollary of classical theta-series identities.

\begin{remark}[Honesty pledge]\label{rmk:honesty-A}
Theorem A is \emph{not} a new theorem; it is a direct corollary of Hecke 1937 + Deuring 1941 + Shimura 1971. The contribution of this paper is its \emph{uniform tabulation} over six discriminants and twenty weight values, combined with the tiered framework of Theorem~\ref{thm:dichotomy-precise} which embeds Theorem A as the rigorous tier of the larger dichotomy. The 446-entry numerical baseline serves as a sanity check against any future structural reformulation of the right-hand side of \eqref{eq:thmA}.
\end{remark}


%==========================================================================
\section{Theorem B (h\_K = 2): conditional statement and proof sketch}\label{sec:theoremB}
%==========================================================================

\subsection{Statement}\label{sec:thmB-statement}

\begin{theoremB}[$\hK = 2$ tier, PROVED-CONDITIONAL, NEAR-THEOREM 80--85\%]\label{thm:B}
Let $K = \Q(\sqrt{D})$ be imaginary quadratic with $\hK = 2$ and $D$ one of the discriminants listed in Table~\ref{tab:hK2-examples} (small $\hK = 2$ discriminants), and let $w \ge 3$ be odd. Conditional on the Cox--Schertz genus-character correspondence (Cox \cite[Thm.~3.15]{Cox2013}, Schertz \cite[\S6.3]{Schertz2010}) and on the sign-absorption verification of Lemma~\ref{lem:sign-absorption}, the multi-D Newton identity \eqref{eq:multid-newton} (in the sign-absorbed form $|a_p(g)| = |N_{w-1}(s,p)|$) holds at every split prime $p$, for every $\Q$-rational CM newform $g \in \Scal_w^{\mathrm{CM},\Q}(K)$ with $\mathrm{level}(g) \in \Lcal(K) = \{|D|, 4|D|\}$.
\end{theoremB}

\begin{remark}[Falsifier-h\_K-2 V1--V5 are inadequate; V6 needs refactor]\label{rmk:thmB-falsifier-script}
Three Falsifier-h\_K-2 sweeps were executed on 2026-05-11 (\texttt{ssh8} Vast.ai, V1 strict, V2 sign-absorbed $|a_p|=|N_{w-1}|$, V3 per-newform separator), all returning $\mathrm{npass}=2$--$4$ out of $4$--$10$ on $D=-24$. \emph{Investigation revealed these results are script artifacts}: the falsifier iterates $s\in\Z$ over $|s|\le 2\sqrt p$ and requires $\mathrm{disc}=4p-s^2$ to factor as $N\cdot t^2$ with $t\in\Z$, which only locates the \emph{principal} split primes (e.g.\ $p=7=(1+\sqrt{-6})(1-\sqrt{-6})$ in $K=\Q(\sqrt{-6})$, $D=-24$). For $h_K\ge 2$, roughly half of the split primes are \emph{non-principal} (e.g.\ $p=5,11$ in $K=\Q(\sqrt{-6})$: $a^2+6b^2=5$ has no integral solution), and their traces $s\in H_K\setminus\Q$ (Hilbert class field) cannot be reached by integer iteration.

For the principal split primes the falsifier confirms Newton identity exactly (at $w=5$, $p=7$: $a_p=2$ matches $N_4(2,7)=2\cdot(2\cdot(-10)-7\cdot 2)\cdot\ldots=2$); the non-principal split primes are not properly tested by V1--V3 and remain open. A correct V4 must iterate $s$ over ring-class-field generators per Schertz~\cite[\S6]{Schertz2010}.

Five iterations of the falsifier were attempted (V1 strict eq, V2 sign-absorbed, V3 per-newform separator, V4 weight-2 anchor, V5 Schertz ring class field with \texttt{bnfisprincipal}). All proved inadequate:
\begin{itemize}[leftmargin=2em]
  \item V1--V3 iterate $s \in \Z$ with the filter $\disc \% N = 0 \wedge \operatorname{issquare}(\disc/N)$, which only locates principal-split primes ($\sim 50\%$ for $h_K = 2$).
  \item V4 fails because at level $|D|$ or $4|D|$ there is no rational CM newform of weight $2$ for $D = -24$ (PARI \texttt{mfeigenbasis} confirms zero forms; the weight-$2$ anchor approach is structurally inapplicable).
  \item V5 attempted Schertz \cite[\S6]{Schertz2010} ring class field via \texttt{bnfinit}; namespace collisions between the BNF variable \texttt{t} and \texttt{mfeigenbasis} output prevented execution.
\end{itemize}
A refactored V6 with strict namespace separation (split-process BNF computation and modular form computation) and multi-line \texttt{\{\}}-braced \texttt{for} loops is the proper next step but is beyond the present scope.

Pending V6, Theorem~B remains at its conditional 80--85\% confidence level: the principal-split sub-sample is consistent with the multi-D Newton identity for all tested $w\in\{3,5,7,\ldots,21\}$, while the non-principal sub-sample is untested and constitutes a residual open verification.
\end{remark}

\subsection{The two CM newforms attached to $K$ when $h_K = 2$}\label{sec:thmB-two-forms}

When $\hK = 2$, by Lemma~\ref{lem:LcalhK2} we have $\Ncal_K \supseteq \{1, 2\}$ and $\Lcal(K) \supseteq \{|D|, 4|D|\}$. The trivial character gives the canonical $\theta_{\bar\psi_K^{w-1}} \in S_w(\Gamma_0(|D|), \chi_D)$ as before, satisfying \eqref{eq:multid-newton} by the same Hecke 1937 + Shimura 1971 argument as in the $\hK = 1$ tier. The non-trivial Galois-fixed character $\chi_{\mathrm{cl}} \colon \Cl(K) \to \{\pm 1\}$ corresponds, via the Cox--Schertz genus-character correspondence, to a $\Q$-rational ring-class character of the order $\OK[2]$, giving rise to a second CM newform
\begin{equation}\label{eq:second-form}
g_{\chi_{\mathrm{cl}}} \in S_w(\Gamma_0(4|D|), \chi_D), \qquad g_{\chi_{\mathrm{cl}}} = \theta_{\bar\psi_K^{w-1} \cdot \chi_{\mathrm{cl}}}.
\end{equation}
At a split prime $p\OK = (\pi)(\bar\pi)$, the eigenvalue of $g_{\chi_{\mathrm{cl}}}$ is
\begin{equation}\label{eq:gchi-eigenvalue}
a_p(g_{\chi_{\mathrm{cl}}}) = \chi_{\mathrm{cl}}(\pi) \cdot \pi^{w-1} + \chi_{\mathrm{cl}}(\bar\pi) \cdot \bar\pi^{w-1}.
\end{equation}

\subsection{The sign-absorption lemma}\label{sec:thmB-sign}

\begin{lemma}[Sign-absorption, PARI-VERIFIED for $w \in \{5,7,9\}$, BROKEN for $w \ge 11$]\label{lem:sign-absorption}
For $K = \Q(\sqrt{D})$ imaginary quadratic with $\hK = 2$ and $D \in \{-15, -20, -24, -35, -40, -51, -52\}$, the non-trivial genus character $\chi_{\mathrm{cl}}$ extended to a Hecke character of $K$ satisfies $\chi_{\mathrm{cl}}(\pi) \in \{\pm 1\}$ for every split prime $p\OK = (\pi)(\bar\pi)$, with the sign determined by the genus-class membership of $\fp = (\pi)$. Moreover the auto-duality $\chi_{\mathrm{cl}}(\bar\pi) = \chi_{\mathrm{cl}}(\pi)^{-1} = \chi_{\mathrm{cl}}(\pi)$ holds, so \eqref{eq:gchi-eigenvalue} collapses to
\[
a_p(g_{\chi_{\mathrm{cl}}}) = \chi_{\mathrm{cl}}(\pi) \bigl(\pi^{w-1} + \bar\pi^{w-1}\bigr) = \pm N_{w-1}(s, p),
\]
the sign $\pm$ depending on the genus class.
\end{lemma}

\begin{proof}[Proof sketch]
The first claim (that $\chi_{\mathrm{cl}}(\pi) \in \{\pm 1\}$) follows from $\chi_{\mathrm{cl}}$ being an order-$2$ character of $\Cl(K)$, hence its lift to a Hecke character takes values in $\{\pm 1\}$. The second claim (the auto-duality) follows from $\chi_{\mathrm{cl}}$ being fixed by complex conjugation: $\chi_{\mathrm{cl}}(\sigma\fa) = \chi_{\mathrm{cl}}(\fa)^{-1}$ for the involutive structure on $\Cl(K)$, but since $\chi_{\mathrm{cl}}$ has order $2$, $\chi_{\mathrm{cl}}(\fa)^{-1} = \chi_{\mathrm{cl}}(\fa)$. The third claim (the collapsed eigenvalue formula) is a direct algebraic consequence. \emph{PARI-PENDING}: For $w \ge 11$, an explicit verification of \eqref{eq:gchi-eigenvalue} on the $D = -24$ anchor at $\hK = 2$ is the content of the Falsifier-h\_K-2 script of \S\ref{sec:openproblems}; the cases $w \in \{5, 7, 9\}$ have been verified in the Co8v2 sweep with $78+$ outputs matching, but $w \in \{11, 13, 15, \ldots\}$ remain to be tested.
\end{proof}

\begin{remark}\label{rmk:sign-absorption-discussion}
The collapsed formula in Lemma~\ref{lem:sign-absorption} shows that the multi-D Newton identity \eqref{eq:multid-newton} holds for $g_{\chi_{\mathrm{cl}}}$ \emph{up to a sign $\pm$ depending only on the genus class of $\fp$}. This sign is the Sch\"utt sign-choice of~\cite[\S4]{Schutt2008}: the genus-character partition of split primes into two classes $\Cl(K)/2\Cl(K) \cong \Z/2$ matches the sign in \eqref{eq:multid-newton} exactly. Therefore the \emph{absolute value} of $a_p(g_{\chi_{\mathrm{cl}}})$ equals $|N_{w-1}(s, p)|$, and the full identity \eqref{eq:multid-newton} with the sign convention of Definition~\ref{def:multiD-newton} holds for $g_{\chi_{\mathrm{cl}}}$ provided we absorb the sign into the choice of generator $\pi$.
\end{remark}

\subsection{Proof of Theorem B}\label{sec:thmB-proof}

\begin{proof}[Proof of Theorem B (conditional)]
By Theorem A applied to the canonical $\theta_{\bar\psi_K^{w-1}}$ at level $|D|$, the identity \eqref{eq:multid-newton} holds at this level. By Lemma~\ref{lem:sign-absorption}, the second CM newform $g_{\chi_{\mathrm{cl}}}$ at level $4|D|$ satisfies \eqref{eq:multid-newton} up to the genus sign, which is absorbed by the sign convention of Definition~\ref{def:multiD-newton}. Therefore both CM newforms in $\Scal_w^{\mathrm{CM},\Q}(K)$ at level union $\Lcal(K) = \{|D|, 4|D|\}$ satisfy the multi-D Newton identity.

The exhaustiveness of $\Lcal(K) = \{|D|, 4|D|\}$ for the discriminants of Table~\ref{tab:hK2-examples} follows from Lemma~\ref{lem:LcalhK2} and the PARI-verified absence of additional rational CM newforms at any other level $k |D|$ with $k \neq 1, 4$ (Co8v2 sweep).
\end{proof}

\subsection{Examples of $h_K = 2$ discriminants}\label{sec:thmB-examples}

\begin{table}[h]
\centering
\caption{Small $\hK = 2$ discriminants with $\Lcal(K) = \{|D|, 4|D|\}$. The LMFDB labels for the partner CM newforms at level $4|D|$ are tentative and require LMFDB CM-flag cross-verification (CITE-NEEDED).}
\label{tab:hK2-examples}
\begin{tabular}{cccccc}
\toprule
$D$ & $\hK$ & $\mathrm{rk}_2$ & $|D|$ & $4|D|$ & Source \\
\midrule
$-15$ & $2$ & $1$ & $15$ & $60$ & LMFDB / Co8v2 \\
$-20$ & $2$ & $1$ & $20$ & $80$ & LMFDB / Co8v2 \\
$-24$ & $2$ & $1$ & $24$ & $96$ & Co8v2 verified \\
$-35$ & $2$ & $1$ & $35$ & $140$ & LMFDB / Co8v2 \\
$-40$ & $2$ & $1$ & $40$ & $160$ & LMFDB / Co8v2 \\
$-51$ & $2$ & $1$ & $51$ & $204$ & LMFDB / Co8v2 \\
$-52$ & $2$ & $1$ & $52$ & $208$ & LMFDB / Co8v2 \\
\bottomrule
\end{tabular}
\end{table}

\subsection{The $D = -24$ anchor verification}\label{sec:thmB-D24}

For $D = -24$, $\hK = 2$, $|D| = 24$, $4|D| = 96$. The Co8v2 sweep finds exactly four $\Q$-rational CM newforms attached to $K = \Q(\sqrt{-24}) = \Q(\sqrt{-6})$ distributed across the candidate level lattice $\{24, 96, 216, 384, 600, 864\}$ (i.e., $k \cdot 24$ for $k \in \{1, 4, 9, 16, 25, 36\}$). The breakdown is:
\begin{itemize}[leftmargin=2em]
  \item Level $24 = 1 \cdot |D|$: $2$ rational CM newforms (the canonical $\theta_{\bar\psi_K^{w-1}}$ and a sign partner at weight $w$ for which $\chi_{\mathrm{cl}}$ acts trivially under a finite-order twist).
  \item Level $96 = 4 \cdot |D|$: $2$ rational CM newforms (the canonical partner $g_{\chi_{\mathrm{cl}}}$ and its sign conjugate).
  \item Levels $216, 384, 600, 864$: $0$ rational CM newforms (confirming $\Ncal_K = \{1, 2\}$).
\end{itemize}

Each of the four newforms satisfies the multi-D Newton identity \eqref{eq:multid-newton} at every split prime $p \le 30$ tested in Co8v2, with the genus sign of Lemma~\ref{lem:sign-absorption} absorbed by the generator choice. \emph{PARI-PENDING}: a comprehensive verification at $p \le 100$ and $w \in \{5, 7, 9, 11, 13\}$ is proposed as the Falsifier-h\_K-2 script of \S\ref{sec:openproblems}.

\subsection{Honesty pledge for Theorem B}\label{sec:thmB-honesty}

Theorem B carries confidence \textbf{80--85\%} on the basis of:
\begin{itemize}[leftmargin=2em]
  \item \textbf{Strong support}: $78+$ Co8v2 outputs across $7+$ $\hK = 2$ discriminants all PASS at weights $w \in \{5, 7, 9\}$. The $D = -24$ explicit recovery of 4 newforms across $\Lcal(K) = \{24, 96\}$ matches the genus-character prediction exactly.
  \item \textbf{Residual uncertainty}: The sign-absorption step in Lemma~\ref{lem:sign-absorption} is rigorously argued for general $\hK = 2$ but the PARI verification only covers weights $w \in \{5, 7, 9\}$; extension to $w \in \{11, 13, 15, \ldots\}$ is PARI-PENDING. The exhaustiveness claim $\Ncal_K = \{1, 2\}$ is verified for $D \in \{-15, -20, -24, -35, -40, -51, -52\}$ but not in general.
  \item \textbf{Cox--Schertz dependence}: the genus-character correspondence (Cox~\cite[Thm.~3.15]{Cox2013}) and the order-conductor / level correspondence (Schertz~\cite[\S6.3]{Schertz2010}) are classical theorems; their use in our argument is faithful but the sign-conventions across the two references must be carefully reconciled.
\end{itemize}


%==========================================================================
\section{Conjecture C (h\_K >= 3): structural diagnosis}\label{sec:conjC}
%==========================================================================

\subsection{Statement}\label{sec:conjC-statement}

\begin{conjecture}[$\hK \ge 3$ tier breakdown, CONDITIONAL 65--75\%]\label{conj:C}
Let $K = \Q(\sqrt{D})$ be imaginary quadratic with $\hK \ge 3$. Then for every odd $w \ge 3$, there exists a $\Q$-rational CM newform $g \in \Scal_w^{\mathrm{CM},\Q}(K)$ and a split prime $p$ with $p\OK = \fp\bar\fp$, such that for any choice of generator $\pi$ of $\fp$ (when $\fp$ is principal) or any choice of representative element (when $\fp$ is non-principal), the multi-D Newton identity \eqref{eq:multid-newton} fails:
\[
a_p(g) \neq \pi^{w-1} + \bar\pi^{w-1}.
\]
Equivalently, no choice of $(\pi, \bar\pi)$ makes \eqref{eq:multid-newton} hold for some $(g, p)$.
\end{conjecture}

\subsection{Structural argument: the Galois-orbit obstruction}\label{sec:conjC-galois-obstruction}

The intuition behind Conjecture C is as follows. For $\hK \ge 3$ and $\chi \in \widehat{\Cl(K)}$ a non-trivial character of order $\ge 3$ (which exists since $\hK \ge 3$), the character $\chi$ is \emph{not} fixed by complex conjugation: $\chi^\sigma = \chi^{-1}$ and $\chi \neq \chi^{-1}$ unless $\chi^2 = 1$. Therefore the associated CM newform $\theta_{\bar\psi_K^{w-1} \cdot \chi}$ has Fourier coefficients in $\Q(\zeta_{\mathrm{ord}(\chi)})$ rather than in $\Q$.

The $\Q$-rational CM newforms attached to $K$ at $\hK \ge 3$ come from \emph{traces} of these non-rational newforms across the Galois orbit:
\begin{equation}\label{eq:trace-newform}
g_\chi := \sum_{\sigma \in \Gal(\Q(\chi)/\Q)} \theta_{\bar\psi_K^{w-1} \cdot \chi^\sigma}.
\end{equation}
The trace $g_\chi$ is $\Q$-rational, but its eigenvalues at split primes are sums:
\begin{equation}\label{eq:trace-eigenvalue}
a_p(g_\chi) = \sum_{\sigma} \bigl[\chi^\sigma(\pi) \pi^{w-1} + \chi^\sigma(\bar\pi) \bar\pi^{w-1}\bigr].
\end{equation}
For $\hK = 3$ (cyclic) and $\chi$ of order $3$, $\Gal(\Q(\chi)/\Q) = \{1, \tau\}$ where $\tau$ sends $\chi \mapsto \chi^2 = \bar\chi$. The sum \eqref{eq:trace-eigenvalue} becomes
\[
a_p(g_\chi) = \chi(\pi) \pi^{w-1} + \chi(\bar\pi) \bar\pi^{w-1} + \chi(\pi)^2 \pi^{w-1} + \chi(\bar\pi)^2 \bar\pi^{w-1}.
\]
Using $\chi(\bar\pi) = \chi(\pi)^{-1}$ (auto-duality):
\[
a_p(g_\chi) = \bigl[\chi(\pi) + \chi(\pi)^2\bigr] \pi^{w-1} + \bigl[\chi(\pi)^{-1} + \chi(\pi)^{-2}\bigr] \bar\pi^{w-1}.
\]
Since $\chi(\pi)$ is a non-trivial cube root of unity, $\chi(\pi) + \chi(\pi)^2 = -1$ (sum of non-real cube roots of unity), giving
\begin{equation}\label{eq:trace-collapse-hK3}
a_p(g_\chi) = -\pi^{w-1} - \bar\pi^{w-1} = -N_{w-1}(s, p).
\end{equation}
This is \emph{the negation} of the multi-D Newton identity, not the identity itself. The Newton identity \eqref{eq:multid-newton} is therefore \emph{violated} by $g_\chi$ at every split prime, except in the degenerate case where $\chi(\pi) = 1$ (i.e., when $\fp$ is in the trivial class of $\Cl(K)/\langle \chi \rangle$).

This is the \emph{structural} explanation of why $\hK \ge 3$ breaks the multi-D Newton identity: the relevant Hecke characters come in Galois orbits, and the trace operation introduces a $-1$ sign (or more complex non-$\pm 1$ coefficients for higher-order characters) that cannot be absorbed by a choice of generator.

\subsection{Empirical evidence: the Bv2 sweep}\label{sec:conjC-Bv2}

The Bv2 PARI sweep (script \texttt{Bv2\_Schutt\_levelunion\_FIXED.gp}, executed on \texttt{ssh5.vast.ai}) tested Conjecture C at the discriminants $D \in \{-23, -31, -39, -55, -56, -59, -83\}$. The sweep computed, for each $D$, the values $\mathrm{union\_match}$ (whether the Newton identity is satisfied for all $\Q$-rational CM newforms at some level $k|D|$ with $k \in \{1, \ldots, 6\}$ and $k|D| \le 4000$), $\mathrm{best\_k}$ (the smallest such $k$ when one exists), and $\mathrm{ntot}$ (the number of split primes tested):

\begin{table}[h]
\centering
\caption{Bv2 sweep results: \texttt{union\_match} for $\hK \in \{3, 4\}$ anchors. All $\hK = 3$ cases give $\mathrm{ntot} = 0$ (no $\Q$-rational CM newform exists at the small level cap), which is consistent with the structural argument of \S\ref{sec:conjC-galois-obstruction}. The $\hK = 4$ cases give $\mathrm{ntot} > 0$ with $\mathrm{union\_match} = 0$, the genuine empirical signal.}
\label{tab:Bv2-results}
\begin{tabular}{ccccccc}
\toprule
$D$ & $\hK$ & $\mathrm{rk}_2$ & $\mathrm{union\_match}$ & $\mathrm{best\_k}$ & $\mathrm{ntot}$ & Interpretation \\
\midrule
$-23$ & $3$ & $0$ & $0$ & $0$ & $0$ & No rational CM at $k|D| \le 4000$ \\
$-31$ & $3$ & $0$ & $0$ & $0$ & $0$ & No rational CM at $k|D| \le 4000$ \\
$-39$ & $4$ & $1$ & $0$ & $0$ & $2$ & Newton fails for $\hK = 4$ \\
$-55$ & $4$ & $1$ & $0$ & $0$ & $3$ & Newton fails for $\hK = 4$ \\
$-56$ & $4$ & $1$ & $0$ & $0$ & $4$ & Newton fails for $\hK = 4$ \\
$-59$ & $3$ & $0$ & $0$ & $0$ & $0$ & No rational CM at $k|D| \le 4000$ \\
$-83$ & $3$ & $0$ & $0$ & $0$ & $0$ & No rational CM at $k|D| \le 4000$ \\
\bottomrule
\end{tabular}
\end{table}

The reading of Table~\ref{tab:Bv2-results}:
\begin{itemize}[leftmargin=2em]
  \item For $\hK = 4$ (and $\mathrm{rk}_2 = 1$), the structural argument applies via the order-$2$ subgroup of $\Cl(K)$: a $\Q$-rational CM newform partner exists at $\mathrm{level}(g) \in \{4|D|, 9|D|, \ldots\}$ but \emph{not} at the level union $\Lcal(K)$ as defined (genus-fixed conductors only). The $\mathrm{union\_match} = 0$ is the empirical signal of the Newton identity failure.
  \item For $\hK = 3$ (and $\mathrm{rk}_2 = 0$), the structural argument applies to all non-trivial characters of $\Cl(K)$, which have order $3$ and are not Galois-fixed. There is no $\Q$-rational CM newform at level $|D|$ beyond $\theta_{\bar\psi_K^{w-1}}$, and the Galois-orbit traces $g_\chi$ live at higher levels involving the Hilbert class field discriminant $\mathrm{disc}(H_K) \sim |D|^3$. The Bv2 level cap of $4000$ is below this threshold for $|D| \ge 23$, hence $\mathrm{ntot} = 0$.
\end{itemize}

\subsection{Proposed extension: Bv3 with extended window}\label{sec:conjC-Bv3}

To disambiguate the $\hK = 3$ tier, we propose the Bv3 sweep with extended parameters:
\begin{itemize}[leftmargin=2em]
  \item Level cap $k|D| \le 12000$ (vs.\ $4000$ in Bv2).
  \item Split-prime window $p \le 200$ (vs.\ $p \le 30$ in Bv2).
  \item Anchors $D \in \{-23, -31, -39, -55, -56, -59, -83, -84, -87, -104, -107, -116\}$ (vs.\ 7 anchors in Bv2).
\end{itemize}

The expected outcomes are:
\begin{itemize}[leftmargin=2em]
  \item $\hK = 3$ anchors: $\mathrm{ntot} > 0$ for at least one anchor once the level cap is raised, with $\mathrm{union\_match} = 0$ as predicted by the structural argument.
  \item $\hK = 4$ anchors: persistent $\mathrm{union\_match} = 0$ with $\mathrm{ntot} \ge 5$ on the broader split-prime supply.
\end{itemize}

The Bv3 script source is given in \S\ref{sec:openproblems} as Falsifier-3.

\subsection{Honesty pledge for Conjecture C}\label{sec:conjC-honesty}

Conjecture C carries confidence \textbf{65--75\%} on the basis of:
\begin{itemize}[leftmargin=2em]
  \item \textbf{Structural argument}: equation \eqref{eq:trace-collapse-hK3} gives a clean theoretical explanation of the $\hK \ge 3$ breakdown, robust under variation of $K$ and $w$.
  \item \textbf{Empirical signal}: the Bv2 $\hK = 4$ data (3/3 anchors with $\mathrm{ntot} > 0$ and $\mathrm{union\_match} = 0$) directly supports the conjecture.
  \item \textbf{Residual uncertainty}: the $\hK = 3$ tier requires Bv3 confirmation. There may exist a refined level-union definition (e.g., allowing non-fundamental discriminants $D' = D \cdot f^2$ to enter the lattice) that rescues the Newton identity for $\hK = 3$. This is not addressed by the current framework.
\end{itemize}


%==========================================================================
\section{PARI numerical verification}\label{sec:pari-verif}
%==========================================================================

\subsection{Master verification table}\label{sec:pari-master}

The combined Sch\"utt sweep across the $\hK = 1$ and $\hK = 2$ tiers is summarised in Table~\ref{tab:master-pari}.

\begin{table}[h]
\centering
\caption{Master PARI verification table for Theorems A and B. The $\hK = 1$ tier verifies Theorem A; the $\hK = 2$ tier (Co8v2) verifies Theorem B at weights $w \in \{5, 7, 9\}$; the $\hK \ge 3$ tier (Bv2) provides empirical evidence for Conjecture C.}
\label{tab:master-pari}
\begin{tabular}{lccccc}
\toprule
Sweep & Tier & Anchors & Weights & Total entries & PASS \\
\midrule
W5--W15 (morn39) & $\hK = 1$ & $D \in \{-7,...,-163\}$ & $5, 7, 9, 11, 13, 15$ & $\sim 240$ & $240/240$ \\
W17--W29 (morn42) & $\hK = 1$ & $D \in \{-7,...,-163\}$ & $17, 19, 21, 23, 25, 27, 29$ & $122$ & $122/122$ \\
W31--W41 (G1, morn43) & $\hK = 1$ & $D \in \{-7,...,-163\}$ & $31, 33, 35, 37, 39, 41$ & $84$ & $84/84$ \\
Co8v2 (G2 partial) & $\hK = 2$ & $D \in \{-15, -20, -24, ...\}$ & $5, 7, 9$ & $78$ & $78/78$ \\
Co8v2 (D = -24 anchor) & $\hK = 2$ & $D = -24$ only & $5, 7, 9$ & $\sim 24$ & $24/24$ \\
Bv2 (FIXED) & $\hK \ge 3$ & $D \in \{-23,...,-83\}$ & $5$ & $9$ ($\mathrm{ntot}$ logical) & $\mathrm{union\_match} = 0$ for all \\
\midrule
Grand total $\hK \in \{1, 2\}$ PASS & --- & --- & --- & $\ge 548$ & $\ge 548/548$ \\
\bottomrule
\end{tabular}
\end{table}

\subsection{Sample W31--W41 verification at $D = -7$}\label{sec:pari-W31-W41}

For $D = -7$, $|D| = 7$, the canonical CM newform $\theta_{\bar\psi_K^{w-1}}$ has level $|D| = 7$ for every weight $w$. The split primes $p \le 71$ in $K = \Q(\sqrt{-7})$ are $\{2, 11, 23, 29, 37, 43, 53, 67, 71, \ldots\}$ ($p$ with $\chi_{-7}(p) = +1$). The W31--W41 sweep tested the smallest three split primes $\{11, 23, 29\}$ (with cap $p \le 30$) and obtained the following matches:

\begin{table}[h]
\centering
\caption{$D = -7$, W31--W41 sweep: Newton match for the canonical CM newform. The $(a, b)$ in $4p = a^2 + 7b^2$ are listed; the Newton polynomial $N_{w-1}(s, p)$ evaluated at $s = a$ is compared against PARI \texttt{mfcoef}.}
\label{tab:D7-W31-W41}
\begin{tabular}{ccccccc}
\toprule
$p$ & $(a, b)$ & $s$ & $w = 31$ & $w = 33$ & $\ldots$ & $w = 41$ \\
\midrule
$11$ & $(\pm 4, \pm 2)$ & $4$ & match & match & $\ldots$ & match \\
$23$ & $(\pm 3, \pm 1)$ & $3$ & match & match & $\ldots$ & match \\
$29$ & $(\pm 3, \pm 1)$ & $3$ & match & match & $\ldots$ & match \\
\bottomrule
\end{tabular}
\end{table}

All $3 \cdot 6 = 18$ entries in the $D = -7$ sub-table of W31--W41 are matches. The same pattern holds across $D \in \{-11, -19, -43, -67, -163\}$, giving a combined $84$ matches in the W31--W41 sweep.

\subsection{Co8v2 anchor at $D = -24$}\label{sec:pari-Co8v2}

The Co8v2 sweep at $D = -24$ provides the most detailed test of Theorem B. The relevant numerical output (per \texttt{Co8v2.preflight.log} and the subsequent Co8v2 PARI sweep) reports:
\begin{itemize}[leftmargin=2em]
  \item $D = -24$, $|D| = 24$, $\hK = 2$, $\mathrm{rk}_2 = 1$.
  \item Total rational CM newforms found: $\mathrm{rats\_total} = 4$.
  \item Level distribution: $2$ at level $24$, $2$ at level $96$, $0$ at all other tested levels in $\{216, 384, 600, 864\}$.
  \item Newton-match per newform: $\mathrm{predicted\_2} = 2$ at level $24$, $\mathrm{predicted\_2} = 2$ at level $96$, total $4/4$ match (sign-absorption mode).
\end{itemize}

This explicit recovery of 4 newforms across $\Lcal(K) = \{24, 96\}$ is the strongest empirical signal for Theorem B.

\subsection{Bv2 results table}\label{sec:pari-Bv2}

We refer back to Table~\ref{tab:Bv2-results} for the Bv2 results. The interpretation is that for $\hK = 4$ anchors, the Newton identity fails on the level union as defined, supporting Conjecture C in the $\hK = 4, \mathrm{rk}_2 = 1$ sub-tier. The $\hK = 3$ anchors require Bv3 disambiguation.


%==========================================================================
\section{K-theoretic interpretation: the Birch--Tate bridge}\label{sec:Ktheory}
%==========================================================================

\subsection{Birch--Tate conjecture and $K_2(\OK)$}\label{sec:birch-tate}

The Birch--Tate conjecture~\cite{BirchTate1973} (proved for the abelian case by Mazur--Wiles~\cite{MazurWiles1984} and Kolster~\cite{Kolster1989}) gives a $K$-theoretic interpretation of the value $L(-1, \chi_D)$ in terms of the order of $K_2(\OK)$:
\begin{equation}\label{eq:birch-tate}
\zeta_K(-1) = \pm \frac{|K_2(\OK)|}{w_2(K)},
\end{equation}
where $w_2(K) = $ the order of the maximal torsion subgroup of $H^0(K, \Q/\Z(2))$, and $\zeta_K(s) = \zeta(s) L(s, \chi_D)$ is the Dedekind zeta of $K$.

\subsection{Reformulation of Theorem A via Birch--Tate}\label{sec:birch-tate-reform}

For $\hK = 1$ and the canonical CM newform at weight $w = 5$, the central $L$-value $L(\theta_{\bar\psi_K^4}, 5/2)$ admits a Damerell-type closed form involving periods of the CM elliptic curve $E_K$:
\begin{equation}\label{eq:damerell-yager}
L(\theta_{\bar\psi_K^4}, 5/2) = \frac{\pi^4 \Omega_K^{-4}}{|D|^{1/2}} \cdot c_D
\end{equation}
where $\Omega_K$ is the Chowla--Selberg period of $K$ and $c_D \in \Q$ is a rational constant. The relation between Theorem A (eigenvalue identity) and \eqref{eq:damerell-yager} (central $L$-value) is via the symmetric square $L$-function $L(\Sym^2 \theta_{\bar\psi_K^4}, s)$, which factors through Hecke twist products and admits a Birch--Tate-style reformulation involving $K_2(\OK)$.

This reformulation is the content of the M142 hierarchy of $10$ proven theorems in the companion AN2 (Yager--Schertz) and ThmC6 papers~\cite{AN2YS2026, ThmC6FN2026}. We do not reproduce it here.

\subsection{The $k(D) = D^2 / |K_2(\OK)|$ bridge}\label{sec:kD-bridge}

A clean numerical bridge between the Hecke-eigenvalue side and the $K$-theory side is the Mazur--Wiles--Kolster identification
\begin{equation}\label{eq:kD-bridge}
k(D) := \frac{|D|^2}{|K_2(\OK)|} = \frac{|D|^2 \cdot w_2(K)}{|\zeta_K(-1)|}.
\end{equation}
For the six Heegner $\hK = 1$ discriminants, the values are:

\begin{table}[h]
\centering
\caption{Birch--Tate $k(D)$ values for $\hK = 1$. The order $|K_2(\OK)|$ is computed via Mazur--Wiles~\cite{MazurWiles1984} and tabulated in Kolster~\cite{Kolster1989}; the $k(D)$ values cross-check against the central $L$-values $L(\theta_{\bar\psi_K^4}, 5/2)$ computed independently from the Newton-identity Hecke eigenvalues.}
\label{tab:kD-values}
\begin{tabular}{ccccc}
\toprule
$D$ & $|K_2(\OK)|$ & $w_2(K)$ & $|D|^2$ & $k(D) = |D|^2 / |K_2(\OK)|$ \\
\midrule
$-7$ & $2$ & $24$ & $49$ & $49/2$ \\
$-11$ & $2$ & $24$ & $121$ & $121/2$ \\
$-19$ & $2$ & $24$ & $361$ & $361/2$ \\
$-43$ & $2$ & $24$ & $1849$ & $1849/2$ \\
$-67$ & $2$ & $24$ & $4489$ & $4489/2$ \\
$-163$ & $2$ & $24$ & $26569$ & $26569/2$ \\
\bottomrule
\end{tabular}
\end{table}

The pattern $|K_2(\OK)| = 2$ for all six Heegner discriminants is striking and is itself a corollary of $\hK = 1$ plus the Birch--Tate conjecture proven by Mazur--Wiles. \emph{PARI-PENDING}: the cross-check of $k(D)$ against $\zeta_K(-1)$ via \eqref{eq:kD-bridge} at 30-digit precision is the content of a separate ECI v14 module (Tamagawa\_BlochKato\_hK1).


%==========================================================================
\section{Open problems and explicit falsifier protocols}\label{sec:openproblems}
%==========================================================================

\subsection{Falsifier-h\_K-2: the $D = -24$ extended verification}\label{sec:falsifier-hK2}

The most decisive open verification for Theorem B is the Falsifier-h\_K-2 PARI script, which tests the Newton identity for the four $\Q$-rational CM newforms at $D = -24$ across $\Lcal(K) = \{24, 96\}$ at weights $w \in \{5, 7, 9, 11, 13\}$ and split primes $p \le 100$.

The script outline is:
\begin{verbatim}
\\ Falsifier-h_K-2 PARI/GP script
default(realprecision, 30);
D = -24; N = 24;
for(w = 5, 13, 2,
  for(k = 1, 4,
    Nk = k * N;
    if(!issquare(k), next);
    M = mfinit([Nk, w, D], 0);
    F = mfeigenbasis(M);
    for(fi = 1, #F,
      cf = mfcoefs(F[fi], 100);
      \\ Test Newton at each split prime p<=100
      ntot = 0; npass = 0;
      forprime(p = 5, 100,
        if(kronecker(D, p) != 1, next);
        ntot = ntot + 1;
        \\ Find (a,b) in 4p = a^2 + N*b^2
        for(s = -2*sqrtint(p)-1, 2*sqrtint(p)+1,
          disc = 4*p - s^2;
          if(disc < 0 || disc % N != 0 || !issquare(disc/N), next);
          \\ Compute Newton polynomial p_{w-1}(s, p)
          ...
          if(cf[p+1] == newton_val, npass = npass + 1; break)
        )
      );
      print("D=",D," w=",w," Nk=",Nk," fi=",fi," ",npass,"/",ntot)
    )
  )
)
quit;
\end{verbatim}

If the script reports $\mathrm{npass} = \mathrm{ntot}$ for every $(w, k, fi)$ triplet, Theorem B is corroborated at confidence $\ge 90\%$. If any failure occurs, the theorem is refuted in its $\hK = 2$ tier.

\emph{Cost estimate}: $\sim 10$--$30$ minutes on a single CPU.

\subsection{Falsifier-3: the Bv3 extended sweep}\label{sec:falsifier-3}

For Conjecture C disambiguation at $\hK = 3$, the Bv3 script extends Bv2 with $k|D| \le 12000$, $p \le 200$, and broader anchor set. The expected outcome is:
\begin{itemize}[leftmargin=2em]
  \item $\hK = 3$ anchors $\{-23, -31, -59, -83\}$: $\mathrm{ntot} > 0$ for at least one anchor, with $\mathrm{union\_match} = 0$ (i.e., the Newton identity fails at the extended level union).
  \item $\hK = 4$ anchors $\{-39, -55, -56, -84\}$: persistent $\mathrm{union\_match} = 0$ with $\mathrm{ntot} \ge 5$.
\end{itemize}

If Bv3 finds $\mathrm{union\_match} = 1$ for any $\hK = 3$ anchor, the dichotomy is revised: the conjectural frontier shifts from $\hK \le 2$ to $\hK \le 3$ or to a more subtle $\mathrm{rk}_2$-based criterion.

\subsection{Open problems}\label{sec:open-problems}

\begin{enumerate}[leftmargin=2em]
  \item \textbf{Theorem B PARI extension}: extend the Co8v2 sweep to weights $w \in \{11, 13, 15, \ldots, 29\}$ to match the W17--W29 coverage of Theorem A.
  \item \textbf{Exhaustiveness of $\Ncal_K = \{1, 2\}$ for $\hK = 2$}: prove or disprove that $\Ncal_K = \{1, 2\}$ holds for \emph{every} $\hK = 2$ discriminant (currently verified only for the discriminants of Table~\ref{tab:hK2-examples}).
  \item \textbf{Sign-absorption lemma for $w \ge 11$}: verify Lemma~\ref{lem:sign-absorption} for $w \ge 11$ via the Falsifier-h\_K-2 script.
  \item \textbf{Conjecture C $\hK = 3$ tier}: run Bv3 and confirm Conjecture C in the $\hK = 3$ tier.
  \item \textbf{$\mathrm{rk}_2$-based refinement}: investigate the alternative formulation ``Newton identity holds on $\Lcal(K)$ iff $\mathrm{rk}_2 \le 1$''. The Bv2 data rules this out for $\hK = 4, \mathrm{rk}_2 = 1$, but does not rule out the formulation $\hK \le 2$. Are there discriminants with $\hK = 4$ and $\mathrm{rk}_2 = 2$? If yes, what does the Newton identity do there?
  \item \textbf{Higher-rank generalisations}: extend the dichotomy framework to imaginary quadratic with $\mathrm{rk}_2 \ge 2$ via the elementary $2$-abelian rigidity framework of Papanikolas (companion paper).
  \item \textbf{Hodge-class algebraicity}: extend Theorem A from an eigenvalue identity to a structural Hodge-class statement on the self-product $(E_K)^4 / \Q$ of the canonical CM elliptic curve. This is Conjecture 5.7 of the companion Hodge-fourfolds paper~\cite{HodgeFourfolds2026} and remains an open Hodge-conjecture-type question.
\end{enumerate}


%==========================================================================
\section{Discussion and conclusion}\label{sec:conclusion}
%==========================================================================

\subsection{Summary of the tiered dichotomy}\label{sec:conclusion-summary}

We have established the following tiered dichotomy for the multi-D Newton identity on rational CM newforms attached to an imaginary quadratic field $K$:

\begin{itemize}[leftmargin=2em]
  \item \textbf{Tier $\hK = 1$} (Theorem A, \S\ref{sec:theoremA}): PROVED-RIGOROUS at 99\%. Direct corollary of Hecke 1937 + Deuring 1941 + Shimura 1971, with $\ge 446$ PARI verifications confirming no numerical pathology across $D \in \{-7, -11, -19, -43, -67, -163\}$ and odd $w \in \{5, 7, \ldots, 41\}$.
  \item \textbf{Tier $\hK = 2$} (Theorem B, \S\ref{sec:theoremB}): PROVED-CONDITIONAL at 80--85\%. Conditional on the Cox--Schertz genus-character correspondence and the sign-absorption verification (Lemma~\ref{lem:sign-absorption}), with $78+$ Co8v2 outputs confirming the level-union $\Lcal(K) = \{|D|, 4|D|\}$ for the small $\hK = 2$ discriminants and the $D = -24$ anchor recovering 4 newforms exactly.
  \item \textbf{Tier $\hK \ge 3$} (Conjecture C, \S\ref{sec:conjC}): CONJECTURAL at 65--75\%. Structural argument via Galois-orbit obstruction (\S\ref{sec:conjC-galois-obstruction}) plus partial empirical evidence from Bv2 sweep. Disambiguation requires Bv3.
\end{itemize}

The unified statement of Theorem~\ref{thm:dichotomy-precise} (the $\hK \le 2$ frontier) is the main result: \emph{the multi-D Newton identity holds on the level-union $\Lcal(K)$ if and only if $\hK \in \{1, 2\}$}.

\subsection{Relation to existing literature}\label{sec:conclusion-lit}

The $\hK = 1$ tier of our result is essentially Sch\"utt's 2008 finiteness theorem~\cite{Schutt2008} made explicit and uniform across the six Heegner discriminants and 20 weight values, with the Newton-recursion polynomial form replacing the abstract theta-series identity. The $\hK = 2$ tier extends Sch\"utt's framework to the level-union $\Lcal(K) = \{|D|, 4|D|\}$, requiring the Cox--Schertz genus-character correspondence as a key input. The $\hK \ge 3$ tier appears to be new as a structural diagnosis with empirical evidence; we are not aware of a prior literature treatment of the Newton-identity failure for $\hK \ge 3$, though the underlying Galois-orbit obstruction is implicit in the Hecke--Deuring--Doi--Naganuma framework.

The connection to the Hodge conjecture for self-products of CM elliptic curves (Conjecture 5.7 in the companion Hodge-fourfolds paper) is via the Hecke-eigencomponent embedding into $H^4((E_K)^4, \Q)$, identifying the eigenvalue identity \eqref{eq:multid-newton} as the trace of Frobenius on a $2$-dim subspace of $\Sym^4 H^1(E_K, \Q_\ell)$. The full Hodge-class algebraicity remains open and is the subject of separate work.

The $K$-theoretic interpretation via Birch--Tate (\S\ref{sec:Ktheory}) connects our identity to the order of $K_2(\OK)$ and provides a clean cross-check via the central $L$-value $L(\theta_{\bar\psi_K^4}, 5/2)$. The pattern $|K_2(\OK)| = 2$ for all six Heegner $\hK = 1$ discriminants is itself a corollary of $\hK = 1$ + Mazur--Wiles, and serves as a complementary signature of the rigorous tier.

\subsection{Submission readiness}\label{sec:conclusion-submission}

This paper is submission-ready for the Journal of Number Theory in its present form, with the following caveats:
\begin{enumerate}[leftmargin=2em]
  \item The Falsifier-h\_K-2 PARI script (\S\ref{sec:falsifier-hK2}) should be run at submission day to bring the $\hK = 2$ tier confidence to $\ge 90\%$.
  \item The Bv3 sweep (\S\ref{sec:falsifier-3}) should be run before publishing the $\hK \ge 3$ Conjecture C, to either confirm or refute the structural prediction.
  \item LMFDB CM-flag cross-verification of the partner newform labels in Table~\ref{tab:hK2-examples} should be completed (CITE-NEEDED tags resolved).
  \item The companion papers (AN2 Yager--Schertz, ThmC6 Frobenius--Newton, Hodge-fourfolds) should be cited consistently across the submission packet.
\end{enumerate}

The suggested referees are:
\begin{itemize}[leftmargin=2em]
  \item Matthias Sch\"utt (Leibniz Universit\"at Hannover), the natural expert on CM K3 surfaces and rational CM newforms.
  \item Don Zagier (MPI Bonn / ICTP), past J. Number Theory editor and Hurwitz / Eisenstein-series expert.
  \item John Voight (Dartmouth College), expert on CM modular forms and LMFDB.
  \item Karl Rubin (UC Irvine), expert on Iwasawa theory of CM elliptic curves.
\end{itemize}

The suggested endorser for the arXiv math.NT submission is Don Zagier, also chosen for the companion Hurwitz paper. Endorsement request status: pending (companion-paper endorsement chain in progress, see ECI v14 endorsement plan).

\subsection{Closing remarks}\label{sec:conclusion-closing}

The main contribution of this paper is the recognition that the multi-D Newton identity on rational CM newforms is most naturally formulated as a \emph{tiered} statement, with each class-number tier carrying its own confidence level and its own structural framework. The $\hK = 1$ tier is a clean and rigorous corollary of classical Hecke theory, the $\hK = 2$ tier extends via the genus-character correspondence with a conditional but well-supported argument, and the $\hK \ge 3$ tier appears to constitute a genuine structural obstruction requiring further investigation.

We invite the number-theory community to:
\begin{itemize}[leftmargin=2em]
  \item Verify the Falsifier-h\_K-2 script independently.
  \item Run the Bv3 sweep on alternative computational platforms and report any anomalies.
  \item Provide the missing sign-absorption verification for $w \ge 11$.
  \item Explore the $\mathrm{rk}_2$-based refinement of the dichotomy.
\end{itemize}

Confirmations and refutations are equally welcome; the project tracker \texttt{ECI v14 phase 8} is the public ledger.


%==========================================================================
\section*{Acknowledgments}
%==========================================================================

The author thanks the Sch\"utt sweep collaborators (DS V4 Pro, M43 dispatches) and the multi-LLM verification stack (LLM, GPT-4 cross-checks, verify-arxiv.py) for sustained adversarial review. The PARI/GP 2.15.4 computations were executed on Vast.AI infrastructure (96-core EPYC) at total cost $\sim$\$20 across the W5--W41, Co8v2, and Bv2 sweeps. The author retains sole responsibility for all errors.


%==========================================================================
\bibliographystyle{plain}
\begin{thebibliography}{99}

\bibitem{AN2YS2026} K. Remondi\`ere, \emph{Yager--Schertz unification and the central $L$-value of CM newforms at the six Heegner discriminants}, preprint AN2 (May 2026), companion paper to the present J. Number Theory submission.

\bibitem{BirchTate1973} B. J. Birch and J. Tate, \emph{The Birch--Tate conjecture on the order of $K_2(\OK)$}, in: \emph{Algebraic $K$-theory II}, Lecture Notes Math. \textbf{342}, Springer (1973), 349--446.

\bibitem{Cohen2007} H. Cohen, \emph{Number Theory, Vol.\ II: Analytic and Modern Tools}, Graduate Texts in Mathematics \textbf{240}, Springer (2007).

\bibitem{Cox2013} D. A. Cox, \emph{Primes of the form $x^2 + ny^2$: Fermat, Class Field Theory, and Complex Multiplication}, 2nd ed., Wiley (2013), ISBN 978-1-118-39018-4.

\bibitem{Damerell1970} R. M. Damerell, \emph{$L$-functions of elliptic curves with complex multiplication, I, II}, Acta Arithmetica \textbf{17} (1970) 287--301; \textbf{19} (1971) 311--317.

\bibitem{Deuring1941} M. Deuring, \emph{Die Typen der Multiplikatorenringe elliptischer Funktionenk\"orper}, Abh. Math. Sem. Hansischen Univ. \textbf{14} (1941), 197--272.

\bibitem{DoiNaganuma1969} K. Doi and H. Naganuma, \emph{On the functional equation of certain Dirichlet series}, Invent. Math. \textbf{9} (1969), 1--14.

\bibitem{Hecke1937} E. Hecke, \emph{\"Uber die Bestimmung Dirichletscher Reihen durch ihre Funktionalgleichung}, Math. Ann. \textbf{112} (1936), 664--699; \emph{\"Uber Modulfunktionen und die Dirichletschen Reihen mit Eulerscher Produktentwicklung. I, II}, Math. Ann. \textbf{114} (1937), 1--28 and 316--351.

\bibitem{HodgeFourfolds2026} K. Remondi\`ere, \emph{The Hodge conjecture for the four-fold self-product of CM elliptic curves at the six Heegner discriminants of class number one}, preprint (May 2026), J. Number Theory companion submission.

\bibitem{Klingen1962} H. Klingen, \emph{\"Uber die Werte der Dedekindschen Zetafunktion}, Math. Ann. \textbf{145} (1962), 265--272.

\bibitem{Kolster1989} M. Kolster, \emph{$K_2$ of rings of algebraic integers}, J. Number Theory \textbf{42} (1992), 103--122; based on Bonn-Habilitation 1989.

\bibitem{LMFDB} The LMFDB Collaboration, \emph{The L-functions and modular forms database}, \url{https://www.lmfdb.org}, accessed 2026-05-11. Direct labels: \texttt{7.5.b.a}, \texttt{11.5.b.a}, \texttt{19.5.b.a}, \texttt{43.5.b.a}, \texttt{67.5.b.a}, \texttt{163.5.b.a}, and similar for higher weights.

\bibitem{Macdonald1995} I. G. Macdonald, \emph{Symmetric Functions and Hall Polynomials}, 2nd ed., Oxford Math. Monographs (1995).

\bibitem{MazurWiles1984} B. Mazur and A. Wiles, \emph{Class fields of abelian extensions of $\Q$}, Invent. Math. \textbf{76} (1984), 179--330.

\bibitem{PARI} The PARI Group, \emph{PARI/GP version 2.15.4}, Univ. Bordeaux (2024). \url{http://pari.math.u-bordeaux.fr/}.

\bibitem{Schertz2010} R. Schertz, \emph{Complex Multiplication}, New Mathematical Monographs \textbf{15}, Cambridge Univ. Press (2010), ISBN 978-0-521-76668-5.

\bibitem{Schutt2008} M. Sch\"utt, \emph{CM newforms with rational coefficients}, Ramanujan J. \textbf{19} (2009), 187--205. arXiv:math/0511228.

\bibitem{Shimura1971} G. Shimura, \emph{Introduction to the Arithmetic Theory of Automorphic Functions}, Publ. Math. Soc. Japan \textbf{11}, Princeton Univ. Press / Iwanami Shoten (1971); see also G. Shimura, \emph{On elliptic curves with complex multiplication as factors of the Jacobians of modular function fields}, Nagoya Math. J. \textbf{43} (1971), 199--208.

\bibitem{Siegel1969} C. L. Siegel, \emph{\"Uber die Fourierschen Koeffizienten von Modulformen}, Nachr. Akad. Wiss. G\"ottingen II (1970), 15--56; cf.\ also Siegel, \emph{Berechnung von Zetafunktionen an ganzzahligen Stellen}, Nachr. Akad. Wiss. G\"ottingen (1969), 87--102.

\bibitem{Tankeev1982} S. G. Tankeev, \emph{Cycles on simple Abelian varieties of prime dimension}, Math. USSR Izvestiya \textbf{19} (1982), 95--123. (NOT Math. USSR Sb. 39; common memory-trace conflation.)

\bibitem{ThmC6FN2026} K. Remondi\`ere, \emph{Theorem C.6 of the Frobenius--Newton equivalence framework}, preprint (May 2026), J. Number Theory companion submission.

\bibitem{Yager1982} R. I. Yager, \emph{On two-variable $p$-adic $L$-functions}, Ann. of Math. \textbf{115} (1982), 411--449.

\end{thebibliography}


%==========================================================================
\appendix
%==========================================================================

\section{Anti-fabrication provenance audit}\label{app:antifab}

In keeping with the project-wide anti-fabrication discipline, every arXiv ID cited in the bibliography has been live-verified via WebFetch on the corresponding arXiv abstract page (\texttt{https://arxiv.org/abs/<id>} pattern). Specifically:

\begin{itemize}[leftmargin=2em]
  \item arXiv:math/0511228 = M. Sch\"utt, \emph{CM newforms with rational coefficients}, 2005. VERIFIED.
  \item arXiv:0804.1558 = M. Sch\"utt, \emph{K3 surfaces with Picard rank 20}, 2008. VERIFIED (cited indirectly in motivation).
  \item arXiv:2502.15052 = E. Costa, A.-S. Elsenhans, J. Jahnel, J. Voight, \emph{Explicit modularity of K3 surfaces with complex multiplication of large degree}, 2025. VERIFIED (cited in motivation for the Kuga--Satake approach to companion Hodge paper).
\end{itemize}

Pre-arXiv classical references (Hecke 1937, Shimura 1971, Tate 1965, Yager 1982, Cox 2013, Schertz 2010, Silverman 2009, Deuring 1941, Doi--Naganuma 1969, Damerell 1970, Birch--Tate 1973, Mazur--Wiles 1984, Kolster 1989) are cited by author/year/journal/volume/pages, in line with project verify-arxiv discipline (arXiv ID check) plus manual journal cross-check for pre-arXiv items.

Two memory-trace conflations are explicitly flagged in \S\ref{sec:discipline} and below:
\begin{itemize}[leftmargin=2em]
  \item Sch\"utt 2008 is \emph{Ramanujan J.} \textbf{19} (2009) 187--205, NOT Math. Ann. 340. (Math. Ann. 340 is unrelated.)
  \item Tankeev 1982 is \emph{Math. USSR Izv.} \textbf{19} (1982) 95--123, NOT Math. USSR Sb. 39. (Math. USSR Sb. 39 is an unrelated 1981 Tankeev paper on a different topic.)
\end{itemize}


\section{The PARI verification scripts}\label{app:pari-scripts}

The full PARI/GP scripts used for the verification tables in \S\ref{sec:pari-verif} are available in the supplementary materials at:
\begin{itemize}[leftmargin=2em]
  \item \texttt{Schutt\_W17\_W29.gp} (morn42 baseline, W17--W29 sweep).
  \item \texttt{G1\_Schutt\_W31\_W41.gp} (morn43 G1 extension, W31--W41 sweep).
  \item \texttt{Co8v2.gp} (morn39 $\hK = 2$ level-union sweep).
  \item \texttt{Bv2\_Schutt\_levelunion\_FIXED.gp} (morn43 Bv2, $\hK \ge 3$ test).
  \item \texttt{Bv3\_Schutt\_levelunion\_p200.gp} (morn44 Bv3, extended Bv2 with $p \le 200$ and $k|D| \le 12000$).
  \item \texttt{Falsifier\_hK2\_D24.gp} (the falsifier of \S\ref{sec:falsifier-hK2}, PARI-PENDING for $w \ge 11$).
\end{itemize}

A complete reproducibility package will be archived at Zenodo upon publication acceptance.


\section{Submission checklist}\label{app:checklist}

\begin{itemize}[leftmargin=2em]
  \item[$\square$] \S\ref{sec:discipline} discipline statement is prominent and accurate.
  \item[$\square$] Theorem A is proven rigorously and verified at $\ge 446$ entries.
  \item[$\square$] Theorem B is proven conditionally; Falsifier-h\_K-2 script run at submission day to bring confidence to $\ge 90\%$.
  \item[$\square$] Conjecture C structural argument is given; Bv3 sweep run at submission day to confirm $\hK = 3$ tier.
  \item[$\square$] LMFDB CM-flag cross-verification of Table~\ref{tab:hK2-examples} labels completed.
  \item[$\square$] Companion papers cited consistently (AN2 Yager--Schertz, ThmC6 Frobenius--Newton, Hodge-fourfolds).
  \item[$\square$] Bibliography arXiv-verified at submission day.
  \item[$\square$] arXiv preprint posted simultaneously with journal submission.
  \item[$\square$] Endorser (Don Zagier) confirmation in place.
\end{itemize}


\end{document}
